\documentclass[12pt]{article}
\usepackage[brazilian]{babel}
\usepackage{hyperref}

\title{Guia de migração do abnTeX2 para o abnTeX3 Community}
\author{Projeto abnTeX3 Community}
\date{2026-08-03}

\begin{document}
\maketitle

\section{Objetivo}

Este guia apresenta uma migração incremental. O módulo
\texttt{abntex3-compat} cobre apenas correspondências diretas e emite avisos;
ele não transforma o abnTeX3 em uma réplica da classe abnTeX2.

\section{Classe e perfil}

Troque a classe legada por \texttt{article}, \texttt{report} ou
\texttt{book}. Em seguida, selecione explicitamente o perfil:

\begin{verbatim}
\documentclass[12pt]{report}
\usepackage[
  perfil=academico,
  tipo-trabalho=tese,
  citacoes=autor-data
]{abntex3}
\end{verbatim}

\section{Metadados}

Reúna os setters legados em uma configuração única. O antigo preâmbulo da
folha de rosto deve ser decomposto, pois natureza, objetivo e área são
metadados independentes no abnTeX3.

\begin{verbatim}
\ABNTEXsetup{
  titulo={Título},
  autoria={Nome da autoria},
  instituicao={Instituição},
  natureza={Tese},
  objetivo={apresentada para obtenção do título de Doutor},
  area={Área de concentração},
  orientacao={Nome da orientação},
  local={Cidade},
  data={2026}
}
\end{verbatim}

\section{Fases e elementos}

Use o fluxo do perfil, como \verb|\ABNTEXacademicoexterna|,
\verb|\ABNTEXacademicopretextual|, \verb|\ABNTEXacademicotextual| e
\verb|\ABNTEXacademicopostextual|. Os elementos são comandos com argumentos
explícitos, documentados no manual do perfil.

\section{Bibliografia}

Remova \texttt{abntex2cite}; configure \texttt{citacoes=autor-data} ou
\texttt{citacoes=numerico}, declare o arquivo com
\verb|\addbibresource{arquivo.bib}| e componha as referências com
\verb|\ABNTEXbibliografia|. Execute \texttt{biber}, não BibTeX. Converta
\verb|\citeonline| para \verb|\textcite| e revise citações parentéticas.

\section{Ponte temporária}

Durante uma conversão gradual, carregue \texttt{abntex3-compat} depois de
\texttt{abntex3}. Resolva os avisos um a um e remova o módulo ao concluir.
A matriz completa e os limites estão em \texttt{docs/compatibilidade.md}.

\end{document}
